\section{Institutional comparison: what can Colombia learn from its neighbours?}
The original benchmark draws on the historical DPI Map \citep{institute_for_innovation_and_public_purpose_ucl_dpi_2025}, operator and policy documents, expert interviews and Colombia's own findings. We keep the same five institutional cases and separate domestic system design from what is known about use by migrants and other transnational users. The technical supplement identifies the authorities, dates and units behind the comparison. An unverified feature is not evidence of poor performance or of absence.

Brazil's Pix illustrates a public operator's ability to set common rules and operate settlement and an alias directory \citep{bcbpix}. Peru offers a different example: the central bank phased interoperability in, starting by connecting Yape and Plin rather than asking users to move to a single national wallet \citep{bcrp2025}. Costa Rica's SINPE Móvil illustrates a simple phone-based transfer through participating institutions, but the number is linked to a local colón-denominated account \citep{bccrsinpe}. In Mexico, SPEI handles interbank transfers \citep{banxicospei}, while CoDi provides a mobile payment and collection interface \citep{banxicocodi}. Access to the underlying rail and use of a particular interface are different measures.

At the time of the survey, Colombia's Bre-B linked participating providers for domestic immediate payments. The comparisons help frame choices about operator responsibilities, phased connections, familiar interfaces and account requirements. They do not rank the five countries by migrant inclusion, verify acceptance of specific documents, or establish the effectiveness of cross-border protection. Any comparison of actual uptake, fees, recovery or international delivery needs its own population, period and source.
